\documentclass[11pt]{letter}
\usepackage[T1]{fontenc}
\usepackage[utf8]{inputenc}
\usepackage{lmodern}
\usepackage[margin=1in]{geometry}
\usepackage[hidelinks]{hyperref}
\signature{Miguel Esteban P\'erez Ibarra}
\address{Miguel Esteban P\'erez Ibarra\\
  Escuela Superior de F\'isica y Matem\'aticas\\
  Instituto Polit\'ecnico Nacional, Mexico City, Mexico\\
  \texttt{supersluky@gmail.com}}
\date{\today}
\begin{document}
\begin{letter}{The Editor\\ \textit{The American Mathematical Monthly}}
\opening{Dear Editor,}

I am pleased to submit the manuscript `A Pythagoras Tree Grown from a
Dodecahedron' for consideration as a Research Article in \textit{The American
Mathematical Monthly}.

Bosman's Pythagoras tree is a favorite picture of recreational mathematics,
and its three-dimensional versions are built from cubes.  The article
describes a Pythagoras tree built instead from regular dodecahedra.  It began
as a piece of computer art, and the article tells how its recipe was recovered
from the original files, why it never collides with itself, and how a
classical fact about the six face axes of the dodecahedron, a `Pythagorean
theorem' for the dodecahedron, picks out a special tree in which every child
rests flat on its parent, surface area is conserved from one generation to
the next, and the limit set has dimension exactly two.  It closes with a
short counting argument that decides when the tree can grow from all twelve
faces at once.

I believe the article suits the \textit{Monthly} because it brings together
ideas a broad readership knows, namely the Pythagorean theorem, the Platonic
solids, symmetry groups, and self-similar fractals, around a concrete object
that is easy to picture.  The prerequisites are undergraduate linear algebra
and a little group theory, and the figures carry much of the story.  One
proof is computer-assisted; it reduces an infinite question to finitely many
linear programs, and all code and data are publicly available.

The manuscript is original, has not been published, and is not under
consideration elsewhere.  English is not my native language, and generative
AI was used to help write parts of the text and to help produce the Blender
renders; this is described in the declaration included in the manuscript.  The figures are in color; color
is needed only in the online version, and black-and-white print is
acceptable.  No funding was received, and I have no competing interests to
declare.

\closing{Sincerely,}
\end{letter}
\end{document}
